\chapter{Metodologia}
\label{ch:metodologia}

\section{Introduzione}

In questo capitolo viene presentata la metodologia adottata per lo sviluppo dell'architettura Zero Trust proposta...

\section{Architettura del Sistema}

\myfig{0.8}{architettura_sistema.png}{Schema dell'architettura del sistema proposto}{architettura}

Come illustrato in Figura~\ref{fig:architettura}, il sistema si compone di tre livelli principali...

\subsection{Livello di Autenticazione}

Il primo livello implementa meccanismi di autenticazione multi-fattore basati su...

% TODO: Completare sezione autenticazione

\section{Implementazione del Modello Zero Trust}

\begin{enumerate}
    \item \textbf{Verify Explicitly}: ogni richiesta viene autenticata...
    \item \textbf{Use Least Privilege}: accesso minimo necessario...
    \item \textbf{Assume Breach}: monitoraggio continuo...
\end{enumerate}

\section{Metriche di Valutazione}

Per valutare l'efficacia della soluzione proposta, sono state definite le seguenti metriche:

\mytab{Metriche di valutazione del sistema}{metriche}{
    \begin{tabular}{lcc}
        \toprule
        \textbf{Metrica} & \textbf{Unità} & \textbf{Target} \\
        \midrule
        Latenza di autenticazione & ms & < 100 \\
        Throughput & req/sec & > 1000 \\
        False Positive Rate & \% & < 5 \\
        \bottomrule
    \end{tabular}
}